% Section 3 --- Method. What is fitted, on what, and what each step can rule out.
\section{Method}
\label{sec:method}

Three objects are fitted, and keeping them separate is what makes a negative
result attributable: a \emph{probe} per model, an \emph{alignment map} per pair,
and a \emph{direction} that scores false agreement. Each is fitted on data the
next one is not scored on.

\subsection{The probe}
\label{sec:method:probe}

For a model and a claim we build a contrast pair: the claim followed by
\texttt{Is this claim true? Yes} and the same prefix followed by \texttt{No}. The
two halves differ in a single token, which is deliberate --- the difference
between them is then a statement about the claim rather than about the surface
form. We read the final layer's mean-pooled hidden state for each half. The final layer
is where cross-model linear compatibility has been localized
\citep{gorbett2026ventriloquist}, which is the site this method needs; it also
means a transported probe reads a \emph{late-stage} representation, and we use
that term rather than claiming access to belief.

We compare three linear probes on identical activations, differing only in how
the direction is chosen:

\begin{itemize}
\itemsep0em
\item \textbf{CCS} \citep{burns2023ccs}, unsupervised, fitted by the
  consistency-plus-confidence objective.
\item \textbf{Mass-mean}, the difference of class means, which
  \citet{marks2024geometry} show is the causally faithful direction.
\item \textbf{Logistic regression} on the same features, the standard
  linear-probing baseline.
\end{itemize}

Two implementation points matter enough to state, because getting either wrong
produces a probe that looks healthy and reads nothing. First, CCS requires each
half of the contrast pair to be standardized \emph{separately}. The halves differ
by one token, so their difference contains a large constant direction --- the
token identity, measured at roughly $1.7\times$ the item-to-item variation --- and
a shared normalizer leaves it in place. CCS then drives both loss terms to zero by
reading that direction alone, giving $p^+ \approx 1$ and $p^- \approx 0$ on every
item: perfectly consistent, perfectly confident, and uninformative. Second, and
for the same reason, restarts cannot be ranked by training loss, because the
degenerate solution attains a lower loss ($\sim\!10^{-4}$) than any probe that
tracks truth ($\sim\!10^{-2}$). We reject restarts whose belief output has
near-zero spread across items, which is a label-free criterion, before ranking the
survivors.

\paragraph{Does the probe read truth, or confidence?}
A probe that merely re-encoded the model's stated confidence would be useless for
oversight, since one could read the output instead. We therefore score every
probe on the half of items the model is \emph{most} confident about, against the
model's own signed confidence on the same items. Confidently wrong items are the
ones oversight has to catch, and a probe earns its place only by beating
confidence there.

\subsection{The alignment map}
\label{sec:method:map}

To read one model's probe on another we fit a ridge map between their activation
spaces on paired activations --- the same inputs through both models --- in both
directions. Linearity is the form identifiability results point at
\citep{nielsen2026logit} and the form shown to preserve downstream structure
between independent models \citep{gorbett2026ventriloquist}. Three choices:

\begin{itemize}
\itemsep0em
\item \textbf{Per-dimension standardization}, not mean-centring alone. A handful
  of dimensions carry most of the variance (on one model the largest
  per-dimension standard deviation is $542\times$ the median, and the variance has
  an effective rank near $1$), so an unstandardized least-squares objective, and
  the $R^2$ that scores it, chiefly measure whether the map predicted those few
  dimensions.
\item \textbf{Fitted without labels, on everything available.} The map needs only
  paired inputs, so restricting it to one dataset's labelled split wastes most of
  the data. The fit is sample-limited rather than geometry-limited: held-out
  standardized $R^2$ at the final layer rises from $0.70$ at $3$k rows to $0.83$
  at $16$k and is still climbing.
\item \textbf{Scored by what it carries, not by how well it fits.} We report
  \emph{transfer loss}: the AUROC a probe gives up when read on the other model
  through the map rather than on the model it was fitted on. CKA and map $R^2$ are
  reported beside it as context. They are different questions --- a probe needs one
  direction to survive, those grade all of them --- and they disagree: we observe a
  pair at CKA $0.49$ and $R^2 \approx 0.5$ whose transported probe loses $0.004$.
\end{itemize}

\subsection{The false-agreement direction}
\label{sec:method:direction}

Within the subset where both probes read ``true'', we fit a linear classifier to
separate the genuinely true claims from the false ones. The feature space
concatenates, for both models, the contrast \emph{difference} and the contrast
\emph{mean}, with the second model pushed into the first's space by the map. The
mean is kept deliberately: a confound shared by both halves of a pair survives in
the mean and cancels in the difference, and that confound is what false agreement
is partly made of --- a difference-only space would be engineered not to show it.

Two controls accompany every number. A \textbf{permutation control} refits the
direction on shuffled labels; because the feature space has tens of thousands of
columns, a direction fitted on noise is effectively a random projection and its
AUROC has a wide distribution, so we take several permutations and report the
spread. A margin smaller than that spread is not evidence. An
\textbf{output-confidence baseline} asks whether the activation space is needed
at all.

\subsection{Protocol}
\label{sec:method:protocol}

Splits are group-aware: a true/false pair built from one source question never
straddles the train/test boundary. Within a dataset, probes and map are fitted on
the train split and everything is scored on the test split.

The generalization test is leave-one-dataset-out. For each held-out dataset the
probes and the map are fitted on the other datasets' train splits and the
direction on their test splits, so the direction never learns from items whose
verdicts its own probes could have memorized, and the held-out dataset is unseen
by all three objects. One detail is necessary rather than cosmetic: a probe keeps
its direction across datasets but not its threshold, because activation scale
shifts and the fitted boundary lands in the wrong place. Left uncalibrated this
produced zero agreements on one dataset and $99.7\%$ disagreement on another ---
invisible in AUROC, which is threshold-free, but fatal to a subset defined by a
threshold. We therefore hold each probe's own positive rate fixed at its value
where it was fitted and read the threshold off the deployment scores, which needs
no labels.

\paragraph{Models and data.}
Ten open-weight models spanning four families and $1.7$B to $32$B parameters,
including two base/instruct pairs, giving $49$ ordered pairs in which the overseer
is never larger than the model it oversees. Five datasets: curated true/false
statements \citep{marks2024geometry}, TruthfulQA
\citep{lin2022truthfulqa}, BoolQ, RTE and IMDB. Activations are extracted in
\texttt{fp16} and cached, so every analysis downstream is deterministic and needs
no accelerator.

\paragraph{Reasoning models.}
All but one family are hybrid thinking models. We extract with no chat template
and thinking off, so no reasoning trace is generated. This keeps the comparison
identical across reasoning and non-reasoning models and keeps the output-space
baseline a genuine single-token Yes/No logprob, which it would not be if the next
token began a trace. The probe therefore reads \emph{pre-reasoning} belief, and we
say so rather than calling it belief.

\paragraph{What we do not separate.}
False agreement has two possible causes: the probes faithfully report a belief the
two models share and that belief is wrong, or both probes fail in the same way,
tracking familiarity rather than truth. Distinguishing them requires an
attribution control --- recovering a prominence direction from familiar versus
defamiliarized phrasings of the same claim and testing whether false agreements lie
along it --- which we have not run. Our claims are therefore about the agreement of
probe readouts, and we are explicit where that is weaker than a claim about the
models themselves.
